\section{Kesimpulan}
\label{sec:kesimpulan}

Berdasarkan keseluruhan proses magang dan pengembangan sistem, dapat disimpulkan bahwa proyek "Pengembangan Sistem Monitoring Pendapatan Daerah (MonPD Reborn) Berbasis .NET" telah berhasil diselesaikan dan secara efektif menjawab permasalahan utama yang dihadapi oleh sistem sebelumnya. Proyek ini sukses menghasilkan serangkaian modul fungsional yang tidak hanya menyajikan data secara interaktif, tetapi juga secara fundamental memperbaiki masalah performa sistem yang sebelumnya menjadi kendala operasional utama.

Keberhasilan utama dari proyek ini adalah implementasi arsitektur aplikasi modern menggunakan ASP.NET Core MVC yang menerapkan teknik pemuatan data secara asinkron (AJAX). Secara kuantitatif, sistem MonPD Reborn berhasil mereduksi waktu muat halaman rata-rata dari $\pm120$ detik pada sistem lama menjadi hanya $3-5$ detik, yang merepresentasikan peningkatan performa sebesar 95.8\%. Selain itu, optimasi pengiriman data berhasil menekan besaran \textit{payload} dari satuan megabyte menjadi kurang dari 100 KB per \emph{request}. Dengan pendekatan ini, aplikasi berhasil mengatasi masalah kelambatan karena data hanya dimuat sesuai kebutuhan pengguna dan tidak lagi membebani \textit{server} maupun \textit{browser} saat pertama kali diakses.

Implementasi fitur pemantauan dan analisis yang dikembangkan telah menyediakan alat yang andal bagi Bapenda untuk meningkatkan pengawasan pendapatan daerah. Dampak operasional yang dihasilkan sangat signifikan, dengan estimasi penghematan waktu kerja rata-rata sebesar $\pm57$ menit per hari bagi setiap pegawai pengguna sistem. Dengan demikian, proyek ini tidak hanya menjadi penerapan ilmu teknis rekayasa perangkat lunak, tetapi juga sebuah solusi konkret yang berhasil meningkatkan efisiensi operasional dan performa teknologi informasi di lingkungan kerja Bapenda Kota Surabaya.